\SourceLeaf{160}mache sich aber zum Gesetz, nichts zu verschweigen, damit andere nach Gefallen auch anders
rechnen können. Die Zahlwerthe der Resultate werden, man halte es, wie man wolle, gleiche
Brauchbarkeit haben, aber ihre Zuverlässigkeit riskirt man für zu gross auszugeben, wenn man
mit dem Ausschliessen zu schnell bei der Hand ist. Geschäfte dieser Art scheinen mir schon
mehr Analogie mit dem \textls{Handeln im Leben} zu haben, wo man selten oder nie mathematische
Strenge und Gewissheit hat und wo man sich begnügen muss, nach bester überlegter Einsicht
zu verfahren.

\begin{center}
. . . . . . . . . . . . . . . . .
\end{center}



\Needspace{6\baselineskip}
\begin{center}
\textsc{Gauss} an \textsc{Gerling}. Göttingen, 2. April 1840.
\end{center}

. . . . . . . . . . In Beziehung auf den wissenschaftlichen Inhalt Ihres Briefes will ich bloss auf die
Bedenklichkeit, die Sie Sich, unnöthigerweise, wegen Heterogeneität von Winkel- und
Linien-Messungen machen, ein Paar Worte erwiedern.

Sind $q, q', q'', q'''$ etc. überhaupt durch Beobachtung gefundene Grössen, homogen oder
heterogen ist ganz gleichgültig, und bezeichnet man mit $dq, dq', dq''\ldots$ die
Differenzen zwischen Beobachtung und Rechnung nach irgend einem System von Elementen, so
ist das Princip der Methode der kleinsten Quadrate \textls{durchaus nicht}, dass
\[
dq^2+dq'^2+dq''^2+\text{etc.}
\]
ein Minimum werden soll, sondern dass
$\left(\dfrac{dq}{m}\right)^2+\left(\dfrac{dq'}{m'}\right)^2+\left(\dfrac{dq''}{m''}\right)^2+$etc.
ein Minimum werde, wo $m, m', m''\ldots$ die respectiven bei den einzelnen Daten zu
befürchtenden mittlern Fehler bedeuten. Nur dadurch, dass in der Praxis so überwiegend oft
solche Fälle vorkommen, wo $m=m'=m''=m'''$ etc. gesetzt werden kann oder muss, wird man
verleitet, das erste Énoncé gelten zu lassen und zu vergessen, dass nur das zweite das
allgemein gültige ist.

Um von diesem Princip Gebrauch machen zu können, muss man, wenn auch nicht genau, doch
einigermassen genähert, wenn auch nicht die Grössen
